\chapter{\abstractname}

Fuzzing is a widely used technique for finding bugs, and most bugs found by continuous fuzzing are regressions introduced by recent code changes.
Such regressions are best caught before a change is merged, but fuzzing time is limited at that point, so the fuzzer's success depends heavily on its starting inputs, the seed corpus.
Existing approaches to generating seeds for regression fuzzing consider only the changed code, even though many bugs manifest in code that a change affects but does not modify.
They also generate seeds from scratch, although an existing corpus is often available.

This thesis proposes TACCO (Targeted Affected-Code Corpus Optimization), a pipeline that prepares a targeted seed corpus prior to fuzzing.
It first distills an existing corpus into the seeds most relevant to the code transitively affected by a commit.
A large language model agent then generates additional seeds aimed at affected code that the corpus does not yet cover.
The agent supplies well-formed inputs with the structure needed to reach that code, and the fuzzer mutates them into the malformed inputs that trigger the bug.

A first study on six projects from the BugOSS benchmark compared three starting corpora and two scopes of affected code.
Starting from the distilled corpus and targeting all affected code gave the best trade-off: compared with the original corpus, the resulting corpora increased coverage of affected code on five of six projects, by up to \num{52.3}~percentage points, while being \SIrange{42}{98}{\percent} smaller.
A second study ran \SI{8}{\hour} AFL++ campaigns on twelve projects, six of them not used for tuning.
Compared with the original corpus, the pipeline did not change how often regressions were found, and its median time to exposure was \SI{71}{\minute} shorter, or \SI{31}{\minute} including augmentation time, but this difference was not statistically significant.
These results show that distillation and augmentation yield smaller corpora that cover more affected code and suggest that such corpora can help expose regressions more quickly in time-limited settings.